% ═══════════════════════════════════════════════════════════════
% SPANISCH ARBEITSBLATT: de + artículo (del) + Así se dice
% ═══════════════════════════════════════════════════════════════
% Thema: Besitz ausdrücken + Klassenraumredemittel
% Klasse 7 | Unidad 2 | Mi instituto | DS 03
% ═══════════════════════════════════════════════════════════════

\documentclass[10pt, a4paper]{article}

\usepackage[utf8]{inputenc}
\usepackage[T1]{fontenc}
\usepackage[ngerman]{babel}

% --- SCHRIFTEN ---
\usepackage{palatino}
\usepackage{helvet}
\newcommand{\sanstext}[1]{{\fontfamily{phv}\selectfont #1}}

\usepackage{microtype}
\usepackage[margin=1.4cm, top=1.6cm, bottom=1.3cm]{geometry}
\usepackage{enumitem}
\usepackage{tabularx}
\usepackage{booktabs}
\usepackage{array}
\usepackage{multicol}
\usepackage{tikz}
\usetikzlibrary{calc}

\usepackage{fancyhdr}
\usepackage{lastpage}
\usepackage{needspace}

\usepackage{xcolor}
\usepackage{amssymb}
\usepackage{tcolorbox}
\tcbuselibrary{skins}

% ═══════════════════════════════════════════════════════════════
% FARBEN
% ═══════════════════════════════════════════════════════════════

\definecolor{kopfzeile}{HTML}{1A1A1A}
\definecolor{akzent}{HTML}{C41E3A}          % Spanisch-Karminrot
\definecolor{hellgrau}{HTML}{F2F2F2}
\definecolor{mittelgrau}{HTML}{777777}
\definecolor{dunkelgrau}{HTML}{444444}

% ═══════════════════════════════════════════════════════════════
% VARIABLEN
% ═══════════════════════════════════════════════════════════════

\newcommand{\thema}{de + artículo (del)}
\newcommand{\klasse}{7}
\newcommand{\maxpunkte}{28}

% ═══════════════════════════════════════════════════════════════
% KOPF- UND FUSSZEILE
% ═══════════════════════════════════════════════════════════════

\pagestyle{fancy}
\fancyhf{}
\renewcommand{\headrulewidth}{0pt}
\renewcommand{\footrulewidth}{0pt}
\cfoot{\sanstext{\footnotesize\textcolor{mittelgrau}{\thepage\,/\,\pageref{LastPage}}}}

% ═══════════════════════════════════════════════════════════════
% BOXEN
% ═══════════════════════════════════════════════════════════════

% Grammatik-Box
\newtcolorbox{grammatikbox}{
    colback=hellgrau,
    colframe=hellgrau,
    boxrule=0pt,
    arc=0pt,
    left=8pt, right=8pt, top=8pt, bottom=6pt,
    borderline north={1.5pt}{0pt}{dunkelgrau}
}

% Merk-Box
\newtcolorbox{merkbox}{
    colback=hellgrau,
    colframe=hellgrau,
    boxrule=0pt,
    arc=0pt,
    left=8pt, right=8pt, top=8pt, bottom=6pt,
    borderline north={1.5pt}{0pt}{akzent}
}

% Regel-Tag
\newcommand{\regel}[1]{%
    \tikz[baseline=(X.base)]\node[fill=akzent, text=white,
    font=\sanstext\scriptsize\bfseries, inner sep=2pt, rounded corners=1pt](X){#1};%
}

% ═══════════════════════════════════════════════════════════════
% BEFEHLE
% ═══════════════════════════════════════════════════════════════

% Spanisch: Palatino fett
\newcommand{\es}[1]{\textbf{#1}}

% Deutsch: klein, kursiv, grau
\newcommand{\de}[1]{{\small\textit{\textcolor{mittelgrau}{#1}}}}

% Lücke
\newcommand{\luecke}[1]{\rule{#1}{0.4pt}}

% Punktebox
\newcommand{\punktebox}[1]{\hfill\sanstext{\small[\_\_/#1]}}

% Aufgabentitel
\newcommand{\aufgabe}[2]{%
    \needspace{4\baselineskip}%
    \vspace{10pt}%
    \noindent\sanstext{\textbf{#1}\quad#2}%
    \vspace{5pt}%
    \nopagebreak[4]%
}

% Teilaufgaben
\newlist{teil}{enumerate}{1}
\setlist[teil]{label=\alph*), leftmargin=1.8em, itemsep=4pt, parsep=0pt, topsep=3pt}

\setlength{\parindent}{0pt}
\setlength{\parskip}{4pt}
\setlength{\columnsep}{24pt}

% ═══════════════════════════════════════════════════════════════
% DOKUMENT
% ═══════════════════════════════════════════════════════════════

\begin{document}

% ═══════════════════════════════════════════════════════════════
% HEADER
% ═══════════════════════════════════════════════════════════════

\noindent
\begin{tikzpicture}[remember picture, overlay]
    \fill[kopfzeile] (current page.north west) rectangle +(\paperwidth, -1cm);
    \node[anchor=west, text=white, font=\sanstext\large\bfseries]
        at ([xshift=1.4cm, yshift=-0.5cm]current page.north west) {\thema};
    \node[anchor=east, text=white, font=\sanstext\small]
        at ([xshift=-1.4cm, yshift=-0.5cm]current page.north east)
        {Spanisch\,·\,Klasse \klasse\,·\,Arbeitsblatt};
\end{tikzpicture}

\vspace{0.5cm}

% --- Name/Datum ---
\noindent
\sanstext{\small\textbf{Name:}} \luecke{5cm} \hfill
\sanstext{\small\textbf{Datum:}} \luecke{2.5cm}

\vspace{8pt}

% ═══════════════════════════════════════════════════════════════
% GRAMMATIK-ÜBERSICHT
% ═══════════════════════════════════════════════════════════════

\begin{grammatikbox}
\begin{multicols}{2}

% --- de + artículo ---
{\large\bfseries de + artículo} \de{von + Artikel}

\vspace{3pt}

{\small
\begin{tabular}{@{}lll@{}}
de + el & = & \es{del} \regel{Kontraktion!} \\
de + la & = & de la \\
de + los & = & de los \\
de + las & = & de las \\
\end{tabular}}

\vspace{4pt}

{\footnotesize\textit{Beispiel:} el libro \es{del} alumno \de{(de + el)}}

\columnbreak

% --- Beispiele ---
{\large\bfseries Ejemplos} \de{Beispiele}

\vspace{3pt}

{\small
\begin{tabular}{@{}l@{}}
la silla \es{de la} profesora \\
los libros \es{del} alumno \\
los cuadernos \es{de los} alumnos \\
la mochila \es{de las} alumnas \\
\end{tabular}}

\end{multicols}
\end{grammatikbox}

\vspace{2pt}

% ═══════════════════════════════════════════════════════════════
% AUFGABEN
% ═══════════════════════════════════════════════════════════════

\begin{multicols}{2}

% --- AUFGABE 1 ---
\aufgabe{1}{Ordne zu: Wann benutzt man \es{del}?} \punktebox{4}

{\small Verbinde mit Linien.}

\vspace{6pt}

\begin{tikzpicture}[
    regelbox/.style={draw, rectangle, minimum width=2.4cm, minimum height=0.9cm, font=\small, align=center},
    beispielbox/.style={draw, rectangle, minimum width=3.4cm, minimum height=0.65cm, font=\small}
]
% Linke Spalte: Regeln
\node[regelbox] (R1) at (0,0) {de + el\\[-1pt]{\tiny\textcolor{mittelgrau}{Kontraktion}}};
\node[regelbox] (R2) at (0,-1.1) {de + la\\[-1pt]{\tiny\textcolor{mittelgrau}{keine Kontr.}}};
\node[regelbox] (R3) at (0,-2.2) {de + los\\[-1pt]{\tiny\textcolor{mittelgrau}{keine Kontr.}}};
\node[regelbox] (R4) at (0,-3.3) {de + las\\[-1pt]{\tiny\textcolor{mittelgrau}{keine Kontr.}}};

% Rechte Spalte: Beispiele (gemischt)
\node[beispielbox] (B1) at (4.6,0) {de las alumnas};
\node[beispielbox] (B2) at (4.6,-1.1) {del profesor};
\node[beispielbox] (B3) at (4.6,-2.2) {de la profesora};
\node[beispielbox] (B4) at (4.6,-3.3) {de los alumnos};
\end{tikzpicture}

\vspace{6pt}

% --- AUFGABE 2 ---
\aufgabe{2}{Ergänze \es{del, de la, de los, de las}.} \punktebox{6}

\vspace{4pt}

\begin{teil}
    \item La mochila \luecke{1.8cm} alumno es azul.

    \vspace{2pt}

    \item El bolígrafo \luecke{1.8cm} profesora es rojo.

    \vspace{2pt}

    \item Los libros \luecke{1.8cm} alumnos están en la mesa.

    \vspace{2pt}

    \item La silla \luecke{1.8cm} profesor está delante de la pizarra.

    \vspace{2pt}

    \item El estuche \luecke{1.8cm} alumna está en la mochila.

    \vspace{2pt}

    \item Los cuadernos \luecke{1.8cm} alumnas son nuevos.
\end{teil}

\columnbreak

% --- AUFGABE 3 ---
\aufgabe{3}{Repaso: ¿\es{Hay} o \es{está/están}?} \punktebox{4}

{\small Ergänze die richtige Form.}

\vspace{4pt}

\begin{teil}
    \item ¿Dónde \luecke{1.5cm} el libro del profesor?

    \vspace{2pt}

    \item En la clase \luecke{1.5cm} muchas sillas.

    \vspace{2pt}

    \item La mochila de María \luecke{1.5cm} debajo de la mesa.

    \vspace{2pt}

    \item ¿\luecke{1.5cm} un bolígrafo en tu estuche?
\end{teil}

\vspace{8pt}

% --- AUFGABE 4: Así se dice ---
\aufgabe{4}{Así se dice: Was sagst du?} \punktebox{6}

{\small Wähle das passende Redemittel.}

\vspace{3pt}

{\footnotesize
\fbox{\parbox{0.9\linewidth}{
\es{¿Cómo?} · \es{No comprendo.} · \es{Tengo una pregunta.} \\
\es{¿Cómo se dice ... en español?} · \es{Más despacio, por favor.}
}}}

\vspace{4pt}

\begin{teil}
    \item Der Lehrer spricht zu schnell.

    $\rightarrow$ \luecke{4cm}

    \vspace{2pt}

    \item Du hast etwas nicht verstanden.

    $\rightarrow$ \luecke{4cm}

    \vspace{2pt}

    \item Du weißt nicht, wie man „Radiergummi" auf Spanisch sagt.

    $\rightarrow$ \luecke{4cm}

    \vspace{2pt}

    \item Du hast akustisch nichts verstanden.

    $\rightarrow$ \luecke{4cm}
\end{teil}

\end{multicols}

% ═══════════════════════════════════════════════════════════════
% AUFGABE 5: Volle Breite
% ═══════════════════════════════════════════════════════════════

\vspace{2pt}
\noindent\textcolor{hellgrau}{\rule{\textwidth}{1.5pt}}
\vspace{2pt}

\aufgabe{5}{Beschreibe die Gegenstände.} \punktebox{8}

{\small Schreibe 4 Sätze mit \es{del/de la/de los/de las}. Benutze auch \es{hay} oder \es{estar}!}

\vspace{3pt}

{\footnotesize\textit{Beispiel:} El libro del profesor está encima de la mesa. / Hay muchos cuadernos de los alumnos.}

\vspace{6pt}

1. \luecke{14cm}

\vspace{8pt}

2. \luecke{14cm}

\vspace{8pt}

3. \luecke{14cm}

\vspace{8pt}

4. \luecke{14cm}

% ═══════════════════════════════════════════════════════════════
% MERKKASTEN: Así se dice
% ═══════════════════════════════════════════════════════════════

\vspace{8pt}

\begin{merkbox}
\sanstext{\textbf{Así se dice – Im Unterricht}}

\vspace{3pt}

\begin{multicols}{2}

{\small
\begin{tabular}{@{}ll@{}}
\es{¿Cómo?} & Wie bitte? \\
\es{No comprendo.} & Ich verstehe nicht. \\
\es{Tengo una pregunta.} & Ich habe eine Frage. \\
\end{tabular}}

\columnbreak

{\small
\begin{tabular}{@{}ll@{}}
\es{¿Cómo se dice...?} & Wie sagt man...? \\
\es{Más despacio, por favor.} & Langsamer, bitte. \\
\es{¿Qué significa...?} & Was bedeutet...? \\
\end{tabular}}

\end{multicols}
\end{merkbox}

% ═══════════════════════════════════════════════════════════════
% FOOTER
% ═══════════════════════════════════════════════════════════════

\vfill

\noindent\rule{\textwidth}{0.5pt}

\vspace{4pt}

\hfill\sanstext{\textbf{Punkte:}} \luecke{1.2cm} / \maxpunkte

\end{document}
